% Phase 3 Progress Checkpoint --- Updated Research Proposal
% Compile: pdflatex proposal_updated.tex  (twice, for references)
% This is a NEW file inside Phase_3/. It does not modify, and is not compiled from,
% the frozen Phase_2/proposal/proposal.tex, which remains on file unchanged.
% All quantitative claims are drawn from Phase_3/phase3_progress_report.md and
% Phase_3/experiments/*.json (5-seed re-run) and are audit-verified.
\documentclass[11pt]{article}

\usepackage[
    top=0.8in,
    bottom=0.9in,
    left=0.9in,
    right=0.9in
]{geometry}
\usepackage{times}
\usepackage{amsmath,amssymb}
\usepackage{graphicx}
\usepackage{hyperref}
\usepackage{enumitem}
\usepackage{booktabs}
\usepackage{xcolor}

\setlist{nosep,leftmargin=*}
\hypersetup{colorlinks=true, linkcolor=blue!50!black, citecolor=blue!50!black, urlcolor=blue!50!black}

\title{\textbf{Research Proposal --- Updated for the Phase 3 Progress Checkpoint}\\
Causal Reinforcement Learning with Hidden Confounders and Partial Observability}
\author{Team Members: Soheil Sayah Varg, Mahdi Shirinbayan, Mahdi Ostadmohammadi}
\date{Deep Reinforcement Learning --- Phase 3 Progress Checkpoint}

\begin{document}
\maketitle

\begin{center}
\begin{minipage}{0.92\textwidth}
\small\itshape
\noindent\textbf{Revision note.} This document supersedes nothing on file: our original
Phase 2 proposal remains submitted and unchanged. It is a self-contained, updated version
of that proposal, prepared for the Phase 3 Progress Checkpoint, incorporating a
complete oracle-verified implementation, a corrected 5-seed empirical study, and two
findings that changed since Phase 2 --- including two of our own numbers that we
re-verified and revised downward before this checkpoint (\S\ref{sec:poc},
\S\ref{sec:phase3}). Full logs,
code, and raw results live in \texttt{Phase\_3/}, isolated from the frozen Phase 2
submission.
\end{minipage}
\end{center}

\vspace{4pt}

% =====================================================================
\section{Motivation \& Contribution}

Offline reinforcement learning from logged data is unreliable when the behavior policy
acted on information the learner never sees. In a partially observable MDP (POMDP) whose
\emph{latent} state drives both the logged actions and the outcomes, the data are
\emph{confounded}: a learner that treats the observation as the state systematically
mis-estimates policy value, and no amount of data fixes the bias. This is the central
obstacle to trustworthy offline RL in healthcare, operations, and any domain with
unrecorded drivers of decisions.

Hong, Qi \& Xu (ICML 2024)~\cite{hong2024} give an elegant theory for exactly this
setting: using a pre-decision \emph{negative control} $O_0$, they identify the confounded
POMDP's reward and transition models through two families of \emph{bridge functions}
($b_R$, $b_D$), then plan pessimistically over a policy-independent confidence region. The
paper is, however, \emph{theory only} --- no algorithm, no code, no experiments. Its
neighbors are likewise unrealized empirically for this exact pipeline: the closest
follow-up runs the bridge \emph{estimator} in isolation~\cite{venkatesh2026}, and the
pessimism precedent is theory-only and for the \emph{un}confounded case~\cite{guo2022}.

\textbf{Contribution (capability, not primacy) --- now empirically settled.} We have built
and fully verified, end-to-end, what is to our knowledge the first empirical
implementation that jointly estimates \emph{both} bridge families via the paper's
two-stage procedure and carries them through plug-in value identification to a tractable
pessimistic policy-\emph{selection} step. We deliberately narrowed this from a primacy
claim in our original proposal: we have not reviewed all concurrent model-based
confounded-POMDP work (e.g.\ DPC-MBRL), so we claim only that we can verify ourselves ---
a capability, not a ranking. Per block, the pessimistic inner minimization is the standard
ellipsoid support-function identity (exact, because the stage-2 risk is quadratic and the
plug-in value multilinear) --- \emph{not} new mathematics, and we say so explicitly rather
than implying the closed form itself was the contribution. Our genuinely new content,
confirmed by a complete implementation, is: (i) a \emph{diagnosis} of a failure mode in the
paper's coupled multi-block pessimism, (ii) an honest, quantified account of a
signal-subspace remedy and its limits, and (iii) a full empirical characterization of
\emph{where} the method helps and where it does not, run at two scales with 95\%
confidence intervals.

% =====================================================================
\section{Research Question}

\emph{When the Hong--Qi--Xu confounded-POMDP framework is implemented end-to-end, does
model-based bridge identification actually de-bias policy evaluation and yield good
policies at practical sample sizes, and can its abstract pessimistic $\max$--$\min$ step be
made computationally tractable?} Our Phase 2 proposal anticipated that the answer would be
\textbf{regime-dependent}; Phase 3 confirms this precisely rather than approximately. Our
deliverable is therefore a quantified characterization of \emph{where} de-biasing helps
(strong confounding, exact identification, toy scale: \S\ref{sec:poc}) versus where
estimation variance dominates and the method is beaten by a naive baseline (benchmark
scale, all 15 cells of our confounding-strength $\times$ sample-size sweep:
\S\ref{sec:poc}).

% =====================================================================
\section{Method}
\label{sec:method}

\subsection{Estimation}

Stage 1: conditional mean embedding $\hat\mu_{W_t|x}$ by kernel ridge regression, plus
conditional density/pmf targets $\hat p(y_t|x_t)$. Stage 2: closed-form bridge solve
$\alpha = (G + N_2\lambda_2 I)^{-1}\hat p$, $G = (\Gamma^\top K_W \Gamma) \odot K_{Y''}$,
implemented and verified in \emph{both} the dual (representer) and primal (operator) forms
(agreement $2.4\times10^{-15}$).

\textbf{Regularization (a Phase 2 finding, unchanged).} The paper's RKHS hyperparameter
schedules do \emph{not} transfer to the tabular class: applied literally, they over-shrink
the estimator by $\approx 80\%$ by the third time-step. We grid-calibrated instead: split
$0.7{:}0.3$ ($N_1 \ge N_2$), $\lambda_1 = 1/N_1$, $\lambda_2 = 0.03/\sqrt{N_2}$ --- values
that work, measured rather than derived from the paper.

\textbf{The spectral-monitor fix.} Ill-posedness cannot be monitored by the ordinary
condition number $\mathrm{cond}(G)$ here: because the observation space is structurally
richer than the latent space ($|O|>|S|$), the design matrix carries a null space that pins
$\mathrm{cond}(G)$ artificially flat even as true estimation quality degrades tenfold (we
confirmed this directly with an $O_0$-noise sweep: stage-1 error degrades $10\times$ while
$\mathrm{cond}(G)$ does not move). We instead monitor the smallest \emph{signal}
eigenvalue $\hat\sigma_2$ of the restricted design --- the eigenvalue restricted to the
subspace the data can actually inform. This is the same structural null space responsible
for the pessimism failure mode below, so the fix and the diagnosis share one root cause.

\subsection{Tractable Pessimism, its Failure Mode, and its Repair}
\label{sec:pessimism-method}

Confidence regions are exact ellipsoids centered at $\hat b$; the per-block minimum is
$V(\hat b) - \sqrt{\xi}\,\|g_\pi\|_{H^{-1}}$, $H = \hat T_2 + \lambda_2 I$. The hard part is
the paper's \emph{coupled} multilinear minimization across blocks.

\textbf{Finding, Phase 3: the literal formula is numerically unusable.} Run as written, the
joint multi-block minimization exploits exactly the null-space directions identified
above: because those directions cost the confidence-region geometry almost nothing, the
optimizer walks arbitrarily far into them. At every width that provably contains the truth
($c \ge 0.1$), the pessimistic value diverges to physically impossible magnitudes ---
$-1.15$ at onset, to $-1779$ at $c=10$, against true values of roughly $1.5$--$2.2$ --- and
mis-selects the best policy (regret $0.68$). This is not a tuning artifact: we traced the
mechanism to the interaction between the null-space geometry and the joint optimization,
and it reproduces under a full seed re-run (5-seed check: $-1.15\to-1779$, materially the
same as our original 3-seed numbers, $-1.12\to-1755$).

\textbf{The multilinear joint-minimization repair: Signal-Projected Pessimism.} We restrict
the pessimistic search to the empirically identified signal subspace via an SVD/eigengap
decomposition, using an orthonormal basis $U$ and $H_{\mathrm{sub}} = U^\top H U$. This
keeps $V_{\mathrm{low}}$ informative ($1.35$ at $c=0.03$) and selects the optimal policy
correctly, with \textbf{0.000 selection regret across every seed}.

\textbf{The honest limit (found by our own adversarial self-audit, not by the mentor or a
reviewer).} The repair is a heuristic, not a certified bound. The true bridge functions
leak up to \textbf{21\% of their norm} outside the projected signal subspace, so the
projected confidence region does not actually contain the truth --- its honest joint
coverage is zero. We describe this as a \textbf{coverage/informativeness trade-off}:
vanilla pessimism covers the truth but is uninformative; our projection is informative and
empirically correct but does not provably cover the truth. $V_{\mathrm{low}} \le
V_{\mathrm{true}}$ held in every test we ran, but by empirical sign-alignment, not
certification. We report this as our headline honest finding rather than an overclaimed
fix.

\subsection{Extensions Added During Phase 3 (scope addition, not a backbone change)}
\label{sec:extensions}

In direct response to the benchmark-boundary result (\S\ref{sec:poc}), we implemented two
extensions not in the original proposal, to characterize --- not merely disclose --- the
boundary:

\begin{itemize}
  \item \textbf{Variance-reduction estimator.} A backward-compatible subclass adding
  James--Stein shrinkage and $K$-fold cross-fitting to the benchmark estimator.
  \textbf{Honest result (corrected in our 5-seed re-run, \S\ref{sec:phase3}):} its variance-reduction
  benefit does \emph{not} hold up under a proper 5-seed re-check; we retract an earlier
  claim to the contrary. It gives a modest, noisy $\approx 7\%$ average MAE reduction, not
  a reliable variance-reduction guarantee.
  \item \textbf{Model-free proximal baseline.} An independent implementation of the Shi et
  al.\ (2022) minimax value-bridge estimator~\cite{shi2022} for a genuine head-to-head
  comparison, not a strawman. Two real bugs were caught only by oracle checks during
  development (\S\ref{sec:phase3}); the corrected version is consistent on the toy
  ($0.0076$ max error at $N{=}200\mathrm{k}$) but under-identified by the benchmark's
  binary proxy variable at scale --- a genuine identification limit of that baseline, not a
  bug.
  \item \textbf{Continuous state/action proof of concept.} A toy-scale extension of the
  same bridge-identification and pessimism pipeline to continuous state-action spaces, via
  a kernel/RKHS bridge estimator on a linear-Gaussian environment with an exact
  closed-form oracle; reuses \texttt{ellipsoid\_opt.py}'s pessimism engine unmodified in
  coefficient space. This tests the paper's own stated generality --- its theory is not
  inherently tabular --- rather than assuming it transfers. Full results in
  \S\ref{sec:continuous}; full derivation and code map in
  \texttt{Phase\_3/docs/continuous\_extension.md}.
\end{itemize}

% =====================================================================
\section{Datasets / Environments}
\label{sec:envs}

\textbf{Primary: the public \texttt{clinicalml/gumbel-max-scm} tabular
benchmark}~\cite{oberst2019} (720 observable-core states $\times$ 2 latent states, 8
actions, $T{=}3$). The latent index drives dynamics and behavior but is censored from the
logs; we construct a pre-decision negative control $O_0$ that satisfies the paper's
Assumption 3.1 by construction, plus a per-step noisy marker that guarantees bridge
existence (deterministic masking alone gives rank $720<1440$, i.e.\ no valid one-step
bridge). A confounding-strength knob $\kappa \in \{0, 0.25, 0.5, 0.75, 1.0\}$ lets us sweep
from no confounding to strong confounding. A naive, confounding-blind observation-MDP
evaluator serves as the baseline every method is measured against.

\textbf{Companion: a minimal toy POMDP} ($|S|{=}2, |A|{=}2, |O|{=}3, T{=}3$) with
\emph{exact}, closed-form oracle bridge functions --- the only substrate where recovery
error against ground truth is directly measurable, and where we run our machine-precision
correctness checks (\S\ref{sec:poc}).

\textbf{Continuous companion: a linear-Gaussian confounded POMDP.} Because the paper's
theory is stated for general --- including \emph{continuous} --- spaces (its bridge
functions live in an RKHS), we added a third environment in which state, action,
observation and reward are all real-valued, and the hidden state is observed only through
a noisy proxy $o_t = s_t + \varepsilon_t$ (the classical measurement-error setting proximal
causal inference was built for). Restricting candidate policies to affine, observation-only
maps $a = K o + c$ --- as all our discrete candidate policies already are --- makes the
closed-loop mean trajectory an exact scalar recursion, so this environment too has an
\emph{exact} closed-form oracle for both policy values and bridge coefficients
(\S\ref{sec:continuous}).

Ground truth throughout is exact dynamic programming on the latent SCM. The full pipeline
runs on a laptop in NumPy; no GPU is used or required.

% =====================================================================
\section{Proof of Concept / Empirical Evidence}
\label{sec:poc}

The pipeline is \textbf{already implemented and oracle-verified end-to-end}. Every number
below is machine-generated from \texttt{Phase\_3/experiments/*.json}; the benchmark and toy
pessimism/de-biasing numbers reflect our final 5-seed re-run (up from 2--3 seeds), which we
undertook specifically because two headline claims below turned out to rest on an
unstably small sample of seeds (\S\ref{sec:phase3}).

\subsection{Oracle verification}

\begin{center}
\begin{tabular}{@{}lc@{}}
\toprule
Check & Result \\
\midrule
Policy-value chain (Theorem 3.5) vs.\ exact dynamic programming & $4.4\times10^{-16}$ \\
Estimator vs.\ exact oracle bridges (population limit) & $3.0\times10^{-15}$ \\
Dual vs.\ primal closed-form solve (independent derivations) & $2.4\times10^{-15}$ \\
Analytic pessimism gradients vs.\ finite differences & $4.9\times10^{-10}$ \\
\bottomrule
\end{tabular}
\end{center}

\noindent Every subsequent result is therefore a fact about the \emph{method}, not an
artifact of the code.

\subsection{Bridge recovery (toy; 3 seeds, unchanged since Phase 2 --- see \S\ref{sec:phase3})}

\begin{center}
\begin{tabular}{@{}lcccc@{}}
\toprule
$N$ & 1,000 & 5,000 & 10,000 & 20,000 \\
\midrule
Reward-bridge error & 1.360 & 0.775 & 0.581 & \textbf{0.459} \\
Dynamics-bridge error & 0.732 & 0.422 & 0.312 & \textbf{0.223} \\
\bottomrule
\end{tabular}
\end{center}

\noindent L2 recovery error vs.\ the oracle decreases monotonically at every step
(log--log OLS slopes $N^{-0.36}$, $N^{-0.39}$). This specific check is the one script in
our suite too slow to safely re-run at more seeds before this checkpoint ($\approx$12 minutes for
3 seeds, dominated by an exhaustive-enumeration population check, not the seed loop); it
already persists full per-seed std, and is queued for a 5-seed re-check post-checkpoint.

\subsection{De-biasing (toy, where the method works as intended; 5 seeds)}

The plug-in beats the confounding-blind naive evaluator in \textbf{12/12}
$(N,\text{policy})$ cells, confirmed on 5 seeds. Worst-policy plug-in bias shrinks
$0.104 \to 0.053 \to 0.036$ as $N$ grows from $5\mathrm{k}$ to $50\mathrm{k}$, while naive
bias stays essentially flat at $0.13$--$0.33$ regardless of $N$ --- the textbook signature
of removed \emph{bias} versus irreducible \emph{bias}.

\subsection{Pessimism --- the honest headline (5 seeds)}

\begin{center}
\begin{tabular}{@{}lccccc@{}}
\toprule
Region size $c$ & 0.1 & 0.3 & 1.0 & 3.0 & 10.0 \\
\midrule
Vanilla $V_{\mathrm{low}}$ (best policy) & $-1.15$ & $-10.4$ & $-62.2$ & $-305.2$ &
\textbf{$-1779$} \\
\bottomrule
\end{tabular}
\end{center}

\noindent True policy values lie between roughly $1.5$ and $2.2$. Signal-Projected
Pessimism repairs this to $V_{\mathrm{low}}=1.35$ at $c=0.03$ with \textbf{0.000 selection
regret} (vs.\ vanilla's $0.68$), but leaks $\approx21\%$ of the true bridge norm outside
the projected subspace, so its honest coverage is zero (\S\ref{sec:pessimism-method}).

\subsection{Benchmark boundary --- the most important negative result (5 seeds)}

\begin{center}
\begin{tabular}{@{}lll@{}}
\toprule
Method & MAE range (15-cell sweep) & Behavior with $N$ \\
\midrule
Naive & \textbf{0.013 -- 0.074} & improves with $N$ \\
Model-based (plain) & 0.072 -- 0.230 & beats naive in \textbf{0/15} cells \\
Model-based (shrunk) & 0.063 -- 0.240 & beats naive in \textbf{0/15} cells \\
Model-free proximal~\cite{shi2022} & 0.164 -- 0.180 & \textbf{flat --- does not improve} \\
\bottomrule
\end{tabular}
\end{center}

\begin{center}
\includegraphics[width=0.95\textwidth]{../results/figures/fig6_mb_vs_mf.png}
\end{center}
\begin{center}
\small\textit{Naive vs.\ model-based (plain/shrunk) vs.\ model-free MAE, 5-seed 95\% CIs.
Left: vs.\ confounding strength $\kappa$ at $N{=}20{,}000$. Right: vs.\ sample size $N$ at
$\kappa{=}1$. Naive stays lowest throughout both panels; the wide, overlapping error bars
on MB-plain/MB-js are the visual evidence behind the retracted shrinkage-variance claim
below.}
\end{center}

At the full $720\times2$ scale, across the entire confounding-strength sweep, the
\textbf{naive, confounding-blind baseline beats every bridge-based method in every single
cell}. The de-biasing method's worst performance is at zero confounding
($\kappa=0$), where it is \textbf{2.8--6.8$\times$ worse} than naive --- precisely where
there is nothing to correct. This is the clearest possible signature that the dominant
error source at this scale is \emph{estimation variance}, not the confounding bias the
method targets: inverting many small, noisy per-cell matrices injects more noise than the
bias it removes.

\subsection{Continuous state and action spaces (5 seeds)}
\label{sec:continuous}

The framework is \emph{not} inherently tabular. Re-reading the paper to verify rather than
assume, it states that \emph{``both S and O are continuous,''} that \emph{``function
approximation is inevitable when state and observation spaces are continuous,''} and that
its two-stage procedure is built on kernel methods~\cite{mastouri2021}. Our discrete
implementation is therefore one \emph{instantiation} of the theory, not its full scope. We
tested that generality directly, replacing the tabular estimator with a two-stage
\textbf{kernel/RKHS} bridge estimator (Stage-1 conditional mean embedding with an RBF
instrument kernel; Stage-2 closed-form regularized solve) on the continuous linear-Gaussian
environment of \S\ref{sec:envs}, and reusing our ellipsoidal pessimism machinery
\emph{unmodified} --- it operates purely in coefficient space and never assumed a tabular
origin.

\begin{center}
\begin{tabular}{@{}lc@{}}
\toprule
Check & Result \\
\midrule
Oracle self-consistency (true bridges vs.\ closed-form value) & $1.665\times10^{-16}$ \\
Value-gradient check (analytic/hybrid vs.\ numerical) & $\approx 1\times10^{-11}$ \\
Bridge recovery, reward / dynamics ($N: 300 \to 4{,}000$) & $0.220 \to 0.055$ / $0.240 \to 0.102$ \\
De-biasing vs.\ naive ($N{=}2{,}000$, 95\% CIs) & plug-in wins \textbf{4/4}, $3$--$11\times$ smaller bias \\
Pessimism validity $V_{\mathrm{low}} \le V_{\mathrm{true}}$ (5 widths $\times$ 5 seeds) & \textbf{100\%} \\
\bottomrule
\end{tabular}
\end{center}

\noindent\textbf{A correction to our own framing, from that same re-reading.} The paper's
sentence continues: \emph{``...while the action space A is finite.''} Our environment uses
a real-valued action, so it varies \textbf{two} things at once relative to the paper ---
continuous observations (covered) and continuous actions (not covered). The identification
chain still reproduced the exact oracle to $1.665\times10^{-16}$, so nothing here is
numerically wrong, but this is \textbf{empirical extrapolation beyond the stated theorem,
not an instantiation of it}, and we claim none of the paper's guarantees for it. The
paper-faithful configuration --- continuous $S$/$O$, \emph{finite} $A$ --- is the first
item of our Phase 4 program (\S\ref{sec:fallback}). We flag this rather than relabel the
work, on the same principle that produced the 21\% leakage disclosure above.

\noindent The de-biasing result is the direct continuous analogue of the discrete 12/12
story, with the same bias/variance signature visible in the CIs (naive $-1.255 \pm 0.018$
vs.\ plug-in $-0.401 \pm 0.216$ on the worst policy). \textbf{Honest caveat, held to the
same standard as \S\ref{sec:pessimism-method}: validity is not selection quality.}
$V_{\mathrm{low}} \le V_{\mathrm{true}}$ holds in every test, but \emph{ranking} policies by
$V_{\mathrm{low}}$ is consistently worse than ranking by the plain plug-in value across the
entire width grid tested (regret $0.40$--$0.48$ vs.\ $0.241$), unlike the discrete toy's
$0.000$. Our hypothesis is chain-compounding across the $T{=}3$ stages --- corroborated by
restart gaps growing from $2\times10^{-16}$ to $62.2$ as the region widens --- but it is not
root-caused, so we report it as an open item, not a solved one. Two new bottlenecks also
appear: dense kernel ridge costs $O(N_1^3)$ (capping this proof of concept at $N\approx
4{,}000$ absent a low-rank approximation), and a \emph{finite-dimensional} kernel silently
re-creates the same rank-deficiency pathology diagnosed in
\S\ref{sec:pessimism-method} --- caught only by checking the Gram matrix's rank directly,
since recovery numbers looked correct either way.

% =====================================================================
\section{Phase 3 Progress Update}
\label{sec:phase3}

\subsection{Completed Milestones Since the Phase 2 Proposal}

\begin{itemize}
  \item \textbf{Core pipeline.} The toy and $720\times2$ benchmark environments, both with
  exact oracle bridges/values, and the naive baseline evaluator (\S\ref{sec:envs}).
  \item \textbf{Estimation layer.} The two-stage bridge estimator in both dual and primal
  forms (verified algebraically identical), plus a factored, benchmark-scale version
  avoiding a combinatorial blow-up in the conditioning-set size (\S\ref{sec:difficulties}).
  \item \textbf{Identification and pessimism layer.} The exact Theorem-3.5
  tensor-contraction chain, the closed-form ellipsoidal pessimism solver, and Signal-
  Projected Pessimism (\S\ref{sec:pessimism-method}).
  \item \textbf{Benchmark-scale extensions.} Variance-reduction (James--Stein) estimator
  and an independent model-free proximal baseline (\S\ref{sec:extensions}), plus a
  comprehensive $5\,\kappa \times 3\,N \times 5\,\text{seeds} \times 4\,\text{methods}$
  sweep with 95\% confidence intervals.
  \item \textbf{Continuous state/action extension.} A toy-scale kernel/RKHS proof of
  concept carrying the same identification and pessimism machinery into continuous
  spaces, oracle-verified to $1.665\times10^{-16}$, with 4/4 de-biasing wins and 100\%
  pessimism validity (\S\ref{sec:continuous}).
  \item \textbf{Validation infrastructure.} Four independent regression suites
  (environment/oracle, estimator, plug-in/pessimism, benchmark), all currently green,
  under three rounds of adversarial self-audit that materially changed how we report two
  of our own headline claims.
\end{itemize}

\subsection{Preliminary Findings}

The full numbers are in \S\ref{sec:poc}. In one sentence: the method works exactly as
theorized in the clean, small-scale regime (12/12 de-biasing wins, machine-precision
correctness) and is well-diagnosed but not fully repairable in its pessimism step (a
genuine null-space failure mode, a working but uncertified heuristic remedy), and is beaten
by a naive baseline at realistic scale because estimation variance there outweighs the bias
it targets. \textbf{Our final rigor pass changed two of these findings materially}: the
$\kappa{=}0$ harm ratio was $23$--$46\times$ on 3 seeds and is $2.8$--$6.8\times$ on 5; a
claimed ``shrinkage cuts variance sixfold'' finding does not replicate at all (lower
variance in only 7 of 15 sweep cells) and is retracted rather than carried into the final
paper. Neither was a coding bug --- both were correctly computed statistics over too few
seeds to be stable, which is itself the finding we report in \S\ref{sec:difficulties}.

\subsection{Difficulties Encountered, and How We Resolved Them}
\label{sec:difficulties}

\begin{enumerate}
  \item \textbf{RKHS regularization theory does not transfer to a tabular
  implementation}, over-shrinking the estimator $\approx80\%$ by the third time-step.
  \emph{Resolved} by an empirical grid search (\S\ref{sec:method}).
  \item \textbf{The standard condition-number diagnostic is misleading} here, pinned flat
  by a structural null space regardless of true estimation quality. \emph{Resolved} by
  monitoring the restricted-spectrum signal eigenvalue $\hat\sigma_2$ instead
  (\S\ref{sec:method}).
  \item \textbf{Diagnosing the pessimism blow-up} required tracing the interaction between
  the null-space geometry and the joint multi-block optimization --- not visible from the
  paper, only from watching the number diverge and working backward. \emph{Resolved} by
  the diagnosis itself (\S\ref{sec:pessimism-method}), which motivated the repair.
  \item \textbf{Catching our own overclaim.} Our first framing of Signal-Projected
  Pessimism was ``the repair''; only a subsequent adversarial self-audit surfaced the 21\%
  leakage that invalidates the coverage claim. \emph{Resolved} by reframing it as a
  disclosed heuristic (\S\ref{sec:pessimism-method}).
  \item \textbf{The benchmark variance wall.} Scaling to 720 observations exposed that
  per-cell estimation variance can dominate the structural bias entirely, the opposite of
  the toy story. \emph{Resolved} by characterizing it with a full confounding-strength
  sweep rather than dismissing it as a tuning issue (\S\ref{sec:poc}).
  \item \textbf{Two real bugs in the model-free baseline, caught only by oracle checks.}
  An off-policy bias (true value $1.865$ evaluated to $1.61$, error not shrinking with
  $N$) was fixed with a proper backward recursion; a subsequent $\sim$11,500$\times$11,500
  dense-matrix blow-up ($\approx$2 GB per fit) was fixed by exploiting a block-diagonal
  structure the naive implementation missed. Neither was visible from the derivation alone.
  \item \textbf{Combinatorial blow-up in the benchmark's conditioning set} (billions of
  cells under a direct implementation). \emph{Resolved} by three structural
  factorizations specific to the environment's observation design (hashed encoding,
  core-diagonal bridge support, closed-form $2\times2$ local inversions).
  \item \textbf{Balancing rigor against the timeline.} Every finding above came from
  deliberately adversarial, repeated self-review, which is expensive; \S\ref{sec:fallback}
  states how we are prioritizing the time remaining before Phase 4.
  \item \textbf{Our own 3-seed statistics were not stable enough to report as claims.}
  Bumping the benchmark sweep to 5 seeds changed the $\kappa{=}0$ harm ratio and
  fully retracted the shrinkage-variance claim (\S\ref{sec:phase3}, above). The lesson: a
  claim derived from a \emph{standard deviation}, not just a mean, needs more seeds than a
  claim about the mean alone.
\end{enumerate}

\subsection{Adjustments to the Original Proposal}

\begin{itemize}
  \item \textbf{Novelty claim narrowed} from an implicit primacy claim to a capability
  claim we can verify ourselves (\S1).
  \item \textbf{The pessimism ``reduction'' is now correctly attributed} to the standard
  ellipsoid support-function identity; our new content is the diagnosis and the
  coverage/informativeness finding, not the closed form itself (\S\ref{sec:pessimism-method}).
  \item \textbf{Signal-Projected Pessimism reframed} from ``repair'' to ``heuristic with a
  measured price'' after our own adversarial audit surfaced the 21\% leak
  (\S\ref{sec:pessimism-method}).
  \item \textbf{The benchmark result reframed} from a soft boundary condition to an
  explicit, quantified negative result, generalized from two confounding levels to the
  full 15-cell sweep (\S\ref{sec:poc}).
  \item \textbf{Scope addition, not a backbone change:} the variance-reduction estimator
  and model-free baseline (\S\ref{sec:extensions}) were not in the original proposal; we
  added them to characterize the benchmark boundary rather than leave it unexamined. We
  consider this an extension of the original method, not a change to it, and flag it for
  the mentor's judgment under the course's backbone-change policy.
\end{itemize}

% =====================================================================
\section{Challenges, Fallback Plan, and Roadmap to Phase 4}
\label{sec:fallback}

\subsection{What Might Still Fail, and Our Fallback}

\begin{itemize}
  \item \textbf{Pessimism guarantees.} The coupled inner minimum is heuristic and the
  signal projection provably excludes part of the truth, so $V_{\mathrm{low}} \le
  V_{\mathrm{true}}$ is empirical, not certified. \emph{Fallback:} report the
  heuristic-with-honest-caveat finding as the contribution (lowest risk, already
  complete); a certified alternative (explicitly bounding the excluded-subspace value
  swing) remains a stretch goal only, not a dependency.
  \item \textbf{Estimation variance at scale.} \emph{Confirmed, not merely anticipated}:
  the benchmark result now shows the estimator dominated by naive at every feasible $N$ we
  tested. We tried the identified next step (James--Stein shrinkage / cross-fitting); its
  benefit is real but small and does not close the gap (\S\ref{sec:extensions}).
  \emph{Fallback:} report this as the honest negative result it is, rather than
  characterize it as fully solved.
  \item \textbf{Baselines.} Only a confounding-blind naive baseline and one
  under-identified model-free baseline are in place. \emph{Fallback:} report the current
  comparison, which is already informative and honestly scoped, and note a
  confounding-aware second baseline as a stated limitation rather than a gap we pretend
  isn't there.
  \item \textbf{Further adversarial review of new code} could surface additional issues
  close to the deadline. \emph{Fallback:} our validation suites are fast (minutes) and
  modular, so we re-verify incrementally rather than needing a final big-bang check.
\end{itemize}

\subsection{Roadmap to Phase 4}

We have not pinned an exact Phase 4 date, since our understanding is that the course
timeline has shifted; we intend to confirm it with the mentor. Stages, in priority order:

\begin{enumerate}[label=\textbf{\Alph*.}]
  \item \textbf{Immediate.} Finalize and deliver this checkpoint; no code changes required.
  \item \textbf{Strengthen empirical rigor.} \emph{Done} for the benchmark sweep
  and the toy pessimism/de-biasing checks (5 seeds, 95\% CIs, two claims corrected). The
  toy $N$-scaling sweep remains at 3 seeds (too slow to safely re-run before this checkpoint;
  queued next). Still open: a second, confounding-\emph{aware} baseline beyond naive. The
  continuous extension (\S\ref{sec:extensions}) has its own, separately-tracked open item
  --- diagnosing why pessimistic selection underperforms plug-in selection there ---
  documented in \texttt{docs/continuous\_extension.md} and not gating this roadmap.
  \item \textbf{Decide the pessimism narrative for the final paper.} We recommend
  reporting the heuristic-with-honest-caveat finding as the contribution (lowest risk,
  already complete, fits the course's accepted ``critical analysis'' project form), with a
  certified alternative as a stretch goal only if time permits.
  \item \textbf{Draft the ICML-format final paper.} Nearly all technical content already
  exists in this document and our internal write-ups; remaining work is compression into
  the required structure within the page budget.
  \item \textbf{Mentor-feedback incorporation pass} after comments on the near-final draft.
  \item \textbf{Final polish and presentation prep} --- figure/table cleanup, proofreading,
  page-budget trimming.
  \item \textbf{Continuous-space program.} Re-reading the paper to verify its own scope
  produced a concrete work plan, in priority order: (i) a \textbf{paper-faithful}
  continuous environment --- continuous $S$/$O$ but \emph{finite} $A$, the only
  configuration where Theorem 3.5 and its rates actually apply (\S\ref{sec:continuous});
  (ii) diagnosing the [C4] pessimistic-selection gap; (iii) sweeping the horizon $T$,
  which directly tests the chain-compounding hypothesis rather than merely scaling;
  (iv) more realistic data, subject to a hard constraint --- measuring de-biasing needs
  \emph{ground-truth} policy values, which logged datasets cannot supply, and Assumption
  3.1 needs a valid pre-decision negative control, which they rarely record. Standard
  offline-RL suites fail both tests and are additionally unconfounded by construction, so
  the realistic target is a simulator or semi-synthetic construction in the medical
  setting the paper itself uses as motivation.
\end{enumerate}

% =====================================================================
\begin{thebibliography}{9}
\small
\bibitem{hong2024} M.~Hong, Z.~Qi, Y.~Xu.
\textit{Model-based Reinforcement Learning for Confounded POMDPs}. ICML, 2024.
\bibitem{venkatesh2026} V.~Venkatesh, A.~Malikopoulos.
\textit{Cross-fitted Proximal Learning for Model-Based Reinforcement Learning}. arXiv:2604.05185, 2026.
\bibitem{guo2022} X.~Guo, M.~Cai, Z.~Yang, Z.~Wang.
\textit{Provably Efficient Offline Reinforcement Learning for Partially Observable Markov Decision Processes}.
ICML, 2022.
\bibitem{oberst2019} M.~Oberst, D.~Sontag.
\textit{Counterfactual Off-Policy Evaluation with Gumbel-Max Structural Causal Models}. ICML, 2019.
\bibitem{shi2022} C.~Shi, M.~Uehara, J.~Huang, N.~Jiang.
\textit{A Minimax Learning Approach to Off-Policy Evaluation in Confounded POMDPs}. ICML, 2022.
\bibitem{bennett2024} A.~Bennett, N.~Kallus.
\textit{Proximal Reinforcement Learning: Efficient Off-Policy Evaluation in POMDPs}.
Operations Research, 2024.
\bibitem{mastouri2021} A.~Mastouri et al.
\textit{Proximal Causal Learning with Kernels}. ICML, 2021.
\end{thebibliography}

\end{document}
